\chapter*{Suggested Reading}
\addcontentsline{toc}{chapter}{Suggested Reading}
\markboth{Suggested Reading}{Suggested Reading}

This monograph explores the combinatorial geometry and affine symmetries of uniform composition mosaics on $\mathbb{Z}_{16}$ ($z_4^4$), alongside its foundational square-dimension counterparts on $\mathbb{Z}_4$ and $\mathbb{Z}_9$. The selected works below offer foundational theory in affine automorphism groups, Galois multiplier decimations, distance geometry, and geometric music theory.

\section*{Affine Group Actions and Macro-Orbit Symmetries}

\begin{description}
    \item[\textbf{William Burnside (1911)}]\mbox{}\\
    \emph{Theory of Groups of Finite Order}, 2nd edition, Cambridge University Press, Cambridge.\\
    The classical exposition of the Cauchy--Frobenius orbit-counting lemma (Section~118), providing the foundational tool for Burnside cycle index evaluations under group equivalence.

    \item[\textbf{George P\'olya (1937)}]\mbox{}\\
    \emph{Kombinatorische Anzahlbestimmungen f\"ur Gruppen, Graphen und chemische Verbindungen}, \emph{Acta Mathematica}, 68(1):145--254.\\
    The fundamental theorem on weighted orbit enumeration and cycle index substitutions under permutation group actions.

    \item[\textbf{Harald Fripertinger (1997)}]\mbox{}\\
    \emph{Enumeration of isometry classes of mosaics}, \emph{Discrete Mathematics}, 176(1--3):87--98.\\
    The foundational algebraic paper formalising the classification of cyclic set partitions and mosaic isometry classes via cycle index plethysm.

    \item[\textbf{Harald Fripertinger (1999)}]\mbox{}\\
    \emph{Enumeration of mosaics}, \emph{Discrete Mathematics}, 199(1--3):49--60.\\
    Extends the enumeration of set partition equivalence classes to general group actions, covering affine normalisers and stabiliser decompositions.

    \item[\textbf{Adalbert Kerber (1999)}]\mbox{}\\
    \emph{Applied Finite Group Actions}, 2nd edition, Springer-Verlag, Berlin, Heidelberg.\\
    An exhaustive modern treatise on the constructive generation of orbits, double-coset algorithms, and cycle index polynomials of affine and projective groups.

    \item[\textbf{Emil Artin (1957)}]\mbox{}\\
    \emph{Geometric Algebra}, Interscience Publishers, New York.\\
    The standard algebraic reference on affine transformations, linear automorphisms, and multiplier structures over finite modular rings.
\end{description}

\section*{Metric Energy, Chord Dispersion, and Distance Geometry}

\begin{description}
    \item[\textbf{Isaac J. Schoenberg (1938)}]\mbox{}\\
    \emph{Metric spaces and positive definite functions}, \emph{Transactions of the American Mathematical Society}, 44(3):522--536.\\
    The classical foundation of isometric embeddings into Hilbert spaces and positive definite chordal distance kernels on circular metrics.

    \item[\textbf{Duncan M. Y. Sommerville (1909)}]\mbox{}\\
    \emph{On Certain Periodic Properties of Cyclic Compositions of Numbers}, \emph{Proceedings of the London Mathematical Society}, Second Series, 7:263--313.\\
    Early 20th-century treatise establishing the symmetry properties and periodic recurrence of cyclic integer compositions and chord strides.

    \item[\textbf{Hansraj Gupta (1979)}]\mbox{}\\
    \emph{Enumeration of Incongruent Cyclic $k$-gons}, \emph{Indian Journal of Pure and Applied Mathematics}, 10(8):964--999.\\
    Exact combinatorial classification of incongruent inscribed polygons in cyclic lattices under Euclidean rotation and reflection.
\end{description}

\section*{The Cyclic Lattice Cartography Trilogy}

\begin{description}
    \item[\textbf{Gemma Knee and Antigra Vitti (2026)}]\mbox{}\\
    \emph{The Polygonal Burnsidean: An Atlas of Inscribed Polygons from $\mathbb{Z}_1$ to $\mathbb{Z}_{16}$}, Newcastle upon Tyne.\\
    The companion volume establishing the comprehensive inventory of single inscribed polygon orbits ($k\text{-}i$) to which every constituent mosaic block is relationally anchored.

    \item[\textbf{Gemma Knee and Antigra Vitti (2026)}]\mbox{}\\
    \emph{An Atlas of Composition Mosaics: The Complete Dihedral Universe from $\mathbb{Z}_1$ to $\mathbb{Z}_8$}, Newcastle upon Tyne.\\
    The companion monograph exploring the unconstrained progression of all 509 dihedral composition mosaics across Bell-number partitions from $\mathbb{Z}_1$ to $\mathbb{Z}_8$.
\end{description}

\section*{Geometric Music Theory and Polyrhythmic Hockets}

\begin{description}
    \item[\textbf{Godfried T. Toussaint (2013)}]\mbox{}\\
    \emph{The Geometry of Musical Rhythm: What Makes a ``Good'' Rhythm Good?}, CRC Press, Boca Raton, FL.\\
    A comprehensive investigation into cyclic metric timelines, analysing maximal evenness, chordal dispersion, and Afro-Cuban timeline patterns as inscribed polygons in $\mathbb{Z}_{16}$.

    \item[\textbf{Dmitri Tymoczko (2011)}]\mbox{}\\
    \emph{A Geometry of Music: Harmony and Counterpoint in the Extended Common Practice}, Oxford University Press, New York.\\
    Geometrical formulation of musical scales, pitch-class sets, and voice leading as configurations and trajectories within quotient orbifolds under dihedral and symmetric group actions.

    \item[\textbf{Milton Babbitt (1961)}]\mbox{}\\
    \emph{Set structure as a compositional determinant}, \emph{Journal of Music Theory}, 5(1):72--94.\\
    Foundational study on the partitioning of the modular aggregate into mutually exclusive, affine-related pitch-class and rhythmic blocks.
\end{description}
